\subsection*{QuickSort}

\Alg 
\begin{enumerate}
    \item Выбрать опорный элемент pivot.
    \item Произвести процедуру Partition, то есть распределить элементы в следующем порядке \textbf{(*)}:
    \begin{enumerate}
        \item Элементы $<$ pivot. Положим, что они занимают отрезок $[i_1, i_2].$
        \item Элементы $=$ pivot. Положим, что они занимают отрезок $[i_2 + 1, i_3].$
        \item Элементы $>$ pivot. Положим, что они занимают отрезок $[i_3 + 1, i_4]$
    \end{enumerate}
    \item Вызвать рекурсивно от $[i_1, i_2]$ и $[i_3 + 1, i_4]$.
\end{enumerate}

\subsection*{Пример процедуры Partition}

(Нам дают индексы left, right - начало и конец сортируемого отрезка, pivot - опорный элемент, vec - сортируемый массив) 

\begin{enumerate}
    \item Распределим элементы так, чтобы сначала шли все элементы $<$ pivot, а потом - все остальные:\\
    
    У нас есть индексы r0 и i. Мы поддерживаем инвариант: vec[i] < vec[r0]. Мы также постепенно двигаем i вправо, а r0 влево. Если инвариант перестаёт поддерживаться, то мы меняем местами значения vec[i] и vec[r0].
    \item Теперь распределим элементы так, чтобы выполнялось \textbf{(*)}:\\

    У нас есть индексы l0 и i (Стоит отметить, что l0 изначально стоит в элементы, где vec[l0] $\geq$ pivot. Мы поддерживаем инвариант: vec[l0] < vec[i]. Мы также постепенно двигаем l0 вправо, а i влево. Если инвариант перестаёт поддерживаться, то мы меняем местами значения vec[l0] и vec[i].
\end{enumerate}

Теперь просто найдём индексы, где начинается/заканчивается отрезок из элементов $=$ pivot и вернём пару из таких индексов.

% \begin{lstlisting}
% pair<int, int> Partition(int[] array, int l, int r) {
%   int pivot = ChoosePivot(array, l, r);
%   int i = l, j = r;
%   while(i <= j) {
%     while (array[i] < pivot) { ++i; }
%     while (array[j] > pivot) { --j; }
%     if (i >= j) { break; }
%     Swap(array[i], array[j]); ++i; --j;
%   }
%   return {j, i};
% }
% \end{lstlisting}

\begin{lstlisting}
#include <vector>

std::pair<int, int> Partition(int left, int right, int pivot,
                              std::vector<int>& vec) {
  int r0 = right;
  while (r0 != left && vec[r0] >= pivot) {
    r0--;
  }
  for (int i = left; i < r0; ++i) {
    if (vec[i] >= pivot) {
      while (r0 > i && vec[r0] >= pivot) {
        r0--;
      }
      if (i < r0) {
        std::swap(vec[i], vec[r0]);
        r0--;
      }
    }
  }
  int l0 = left;
  while (vec[l0] < pivot) {
    l0++;
  }
  for (int i = right; i > l0; --i) {
    if (vec[i] == pivot) {
      while (l0 < i && vec[l0] == pivot) {
        l0++;
      }
      std::swap(vec[i], vec[l0]);
      l0++;
    }
  }
  int idx1 = left;
  while (vec[idx1] != pivot) {
    idx1++;
  }
  int idx2 = right;
  while (vec[idx2] != pivot) {
    idx2--;
  }
  return {idx1, idx2};
}
\end{lstlisting}

\Thbd Среднее время работы быстрой сортировки при случайном выборе pivot составит $O(N\,log\,N)$, а доппамять логарифмична.

\pagebreak